\section{伊犁之后：征服不是均质的治理}

1755年清军进入伊犁、击溃达瓦齐；阿睦尔撒纳旋即反清，1758至1759年南路大小和卓政权又告覆亡。战报可以把这些行动排成一条向西推进的胜利线，却不能使天山南北在同一时刻、以同一方式进入清廷的治理。伊犁河谷的驻防、军台、屯田与安置，面对的是牧地、水源和越冬条件；南疆绿洲的伯克、渠水管理者、城镇商人和农户，则在税役、市场与灌溉的安排中面对新政权。将军衙门可以调遣旗人、满汉军人与屯户，也必须依赖地方掌握道路、语言、收获时令和水利的人。征服改变了这些关系，却没有把它们抹平。

供给把这一区别变成日常的政治问题。军粮要沿军台和商路转运，牲畜需有草场，驻军与来迁者又同既有居民争取水、地和贸易机会；一封要求催粮或开垦的奏折，只能证明上级看见了短缺并试图处置，不能证明粮食准时抵达，或某一条渠水的分配已被接受。乌什在1765年的动荡因而不只是驻军体系里的一次“意外”：征发、劳役、地方权威与绿洲生活相互挤压，足以使一座有城、有兵的地方仍然失序。清廷随后能够镇压并重组行政安排，但镇压的完成与社会关系的稳定是两件不同的事。

西藏、青海、尼泊尔边境、川西山地和东南海疆也各有自己的瓶颈。驻藏大臣与噶厦关系的调整、廓尔喀战争后的规条，确实把更多事务接到北京的边务文书中；消息和物资却要穿过寺院网络、山口、驿道和地方政教关系。清缅战争、安南撤兵与台湾林爽文起事所显出的，并非一个简单的“边远地区难以控制”，而是病疫、雨季、航道、翻译、居民合作和外交礼仪会在不同地方以不同方式限制军令。帝国可在数个方向同时用兵，却不能让每一条道路、每一座港口和每一个中介都按中枢的时间表运转。

广州一口通商与马戛尔尼使团把同样的问题推到港口。行商、海关官员、水师、闽粤商人和抵达的欧洲船只并非站在一条命令链上；货物、欠款、引水、停泊与礼仪的规则须在具体交涉中落定。把一口通商写成闭关不动，或把使团写成中国已经“开放”的开端，都会遗漏这些在港口实际承担风险的人。至白莲教战争爆发时，高额军费、山区道路、地方结社与动员已重新缠结起来：扩张并未只留下疆域，也留下需要由户部、地方官、民夫、商人和村寨继续承担的财政与治理压力。

\subsection{材料与争议}

满汉《实录》、军机处档、理藩院档与各类方略最适合追索任命、调兵、赏罚和清廷如何叙述战事；它们让人看见一项命令进入了何种机构，却常把抵达后的协商、拖延与损失压缩为结果。地方奏报同样是官员为考成、请饷或避责而写的陈述。军报中的歼获、伤亡和“平定”因此应当首先读作战争中的国家报告，而不是可以直接累加为人口后果或地方态度的统计。

维吾尔／察合台文、蒙古文和藏文文书，寺院记录、契约与地方档案，可在一些地点补出水利、土地、捐施、贸易和地方机构的时间；俄国、尼泊尔、缅甸及欧洲来往文书又保存了跨境行动者的另一种观察位置。这些材料的留存并不平均。它们不能合成一幅无缝的“地方声音”，正因为如此，缺少自述的社群不应由官方分类或后出的民族名称替其发言。能较稳妥地说的是某项命令被下达、某次战事被报告，或某份契约记录了一项关系；至于服从的范围、迁徙的规模和暴力的总量，仍须按地点、语言和文类分别追问。

\subsection{伊犁、喀什噶尔与绿洲的水}

伊犁将军衙门的文书把驻防、屯田、军台和赏罚编入征服后的秩序。它们可确认清廷试图把兵、粮、土地和人事联结起来，却不能把伊犁河谷的牧地与南疆城镇的渠水当作同一种可调度资源。前者牵连牧群的季节移动与驻防需求，后者牵连绿洲耕作、城镇市场和地方权威；商队也要在关卡、治安和跨境机会之间判断行程。阿睦尔撒纳、大小和卓、伯克、驻防官和商人处于不同的联盟与风险之中，不能被概括为“归附者”与“反叛者”两类。

乌什的情形使这种差别格外清楚。军粮能够进入城中，只说明转运在某一刻成功；它并不回答由谁交纳、谁服役、谁裁断税役与灌溉的纠纷。把《平定准噶尔方略》的凯旋叙述同满文军机档、维吾尔／察合台文契约以及俄国和中亚材料并读，才可在局部重建水、地、贸易与迁徙怎样被重排。即使如此，保存下来的文书也更容易记录官府、商人与能订立契约者，较难保存佃户、妇女或临时迁徙者的选择；不能从档案的沉默推断他们没有承担代价。

\subsection{拉萨、金川与西南山道}

清廷在拉萨调整驻藏大臣及相关规条时，面对的不是一块可由一纸章程均质管理的高原。噶厦、寺院、青海蒙古诸部、廓尔喀使团和驿道上的运输者分别影响消息与供给能否越过山口。藏文寺院文书能够显示某些机构、捐施和地方时间，清廷档案则能够显示驻军、任命与奏议；二者都不能独自代表整个高原，更不能由北京的制度分期替代当地政治的节奏。

川西金川把财政成本压在山道的长度和坡度上。土司关系、碉楼之间的通道、民夫粮运与季节性疾病，使清军即使取得一处据点，也要继续为炮械、粮食和人员的抵达付出代价。所谓改土归流的结果，须由地方档、碑刻、诉讼与迁徙记录追问土地、债务和官署究竟在哪些村寨改变；战报中的“平定”不能预先回答这些问题。金川的战后安排也不能替高原或西南其他地区作结论，正如驻藏体制不能替山地道路解决供给。

\subsection{广州、台湾与海上通路}

广州的行商、海关官员、福建与广东水师、闽南商人与欧洲船只共同限定贸易和防务的边界。一口通商调整了合法贸易的门槛，却没有让海岸社会停止流动：船只的航期、港口的引水、商人的信用和沿岸居民的生计仍使规则在执行中不断被重谈。海关账册和航海日志可以说明一条航线、一批货物或一次交易的制度位置，不能由此推出整个海岸都遵从同一秩序。

1787年林爽文起事的军报记录清军收复城镇，台湾府县志、契约与平码材料则把垦殖、租佃、渡海和地方武装放回具体地点。平原村落、鹿港等港口与内山社群承受的征发和接触，不能被一个“收复”或“设治”概括。原住民社群的自述保存尤其稀少；海防命令、移民契约和殖民或官府记录都带着各自的土地与治理诉求。这里的证据缺口不是可以用更响亮的国家叙事填补的空白。

\section{读者事件簇：乌什的驻防与金川的山道}

\eventref{EQ-1765-USH-REBELLION}{乌什事件（1765）}发生在征服后的南疆驻防体系之内。乾隆三十年的实录、军机处档和战事方略能跟出清廷接到消息、调兵和处理善后的次序；它们的用语也首先服务于镇压与重组治理的需要。若只沿着这条文书链阅读，城内税役、劳役与地方权威的冲突就会被化为抽象的“不服”。维吾尔和察合台文材料、地方文书与跨境商路记录虽各自残缺，却提醒读者，军粮入城和官员到任都不足以证明当地社会已经接受新的权力关系。

\eventref{EQ-1776-JINCHUAN-CONQUEST}{第二次金川战争结束（1776）}把类似的限度放在川西山道上。攻克碉楼结束了一轮军事行动，却没有终止人力和物资的账目：夫役、道路、炮兵以及四川和邻省的转运，已把负担分散给许多并不出现在凯旋叙述中的村寨和家庭。战后设治、土地关系与迁徙的变化需要另以地方材料检验。乌什与金川并非同一场冲突；前者使绿洲驻防的摩擦可见，后者让山地后勤的延续显形。两者共同阻止读者把驻防城或“平定”误读为已经建立了同一种秩序。

\section{读者事件簇：青海部署与高原中介}

\eventref{EQ-1728-ZHUNGHAR-TIBET-SETTLEMENT}{调整驻藏与青海蒙古部署（1728）}，是在拉萨政局和草原力量关系不稳、准噶尔威胁仍在的条件下作出的军政回应。雍正六年的实录、理藩院和相关军政文书可以定位清廷的部署，却不能代替高原上的实施经过。驻防名额写在文书中以后，消息仍要经由青海蒙古诸部、寺院网络、驿道与运输者传递；粮草能否前行，也取决于山口、季节和地方协商。清廷由此更直接地协调西藏事务，并不意味着它预先决定了这些中介会在何处配合、拖延或改写命令。

\section{读者事件簇：乾隆继承的不是空白账册}

\eventref{EQ-1735-QIANLONG-ACCESSION}{乾隆即位（1735）}时，雍正朝的财政整顿、朱批奏折与军机文书已使许多事务能够较快地进入中枢。雍正十三年的实录、奏折与题本能使人看见继承完成和官僚程序的延续，却不能把档案的密集误作地方问题已经被掌握或解决。新皇帝继承的是可以调度的制度，也继承了等待核销的亏空、需要转运的粮饷、山地与边地的供给难题，以及必须由地方官、商人和社群承担的执行成本。此后的扩张正是在这些并未清空的账目上展开。

\section{读者事件簇：第一次金川战争的山道}

\eventref{EQ-1738-FIRST-JINCHUAN-WAR}{大金川战事展开（1738）}时，清廷对川西治理的扩展卷入既有的地方权力冲突。乾隆三年的战报与地方奏报保存了进军、请饷和处置的官方序列；碉楼之间狭窄而曲折的通道、土司间的关系和四川粮运，却使兵力进入山谷不等于控制道路。清军后来以招抚与和议收束战事，显示军事推进必须同地方联盟和供给的限度谈判。不能从一份讲述行动的奏报推出山地社会已经按清廷设想重新排列。

\section{读者事件簇：第二次金川战争的夫役账}

\eventref{EQ-1768-SECOND-JINCHUAN-WAR}{第二次金川战争开始（1768）}，意味着第一次和议未能消除土司与清廷之间的权力冲突。炮兵、粮草和民夫须从四川及邻省沿山道转运；每一次征发或加派都可能暂时接上前线，却把劳力、债务与村寨缺人的风险留在后方。军机处档和战事方略能追踪更大规模资源投入的意图与报告，难以单独计量劳役如何落在不同家庭。战争的代价因而不能只由攻下几座碉楼来计算，它也为1776年之后的设治与迁徙留下了需要继续处理的条件。

\section{读者事件簇：土尔扈特回归的安置难题}

\eventref{EQ-1771-TORGHUT-RETURN}{土尔扈特部回归（1771）}把伏尔加草原、俄国边境与伊犁的牧地安排接到清廷面前。实录、军机处和理藩院档案记录接纳、赏赐与安置的军政语言，也显示这不是单纯的边界新闻：部众迁徙之后，水草、牲畜、疾病、驻防与邻近部落的资源边界都要重新协商。俄国文献、蒙古文材料与清廷档案对迁徙的叙述位置不同，不能由任何一方的数字概括全部损失或内部差异。“安置”不是行政程序的终点，而是草场、军政和跨境关系中一段新的不确定开端。
